\section[局部算子的交换代数]{局部算子的交换代数（4 月 13 日）}

\srcpage{59}

我们想用边界条件（boundary condition）的语言，来解释此前得到的 Hecke 代数（Hecke algebra）故事。先考虑一个 Lagrange 场论（Lagrangian field theory），它是映射 $\Sigma\to Y$ 的理论。这里 $\Sigma$ 是时空（spacetime），$Y$ 称为靶（target）。用量子力学的语言说，我们在把映射空间 $\operatorname{Map}(\Sigma,Y)$ 线性化。

从这个视角看，一类自然的边界条件浮现出来。允许 $\Sigma$ 带边界，并固定一个子流形 $Z\subset Y$，然后只保留满足 $\varphi(\partial\Sigma)\subset Z$ 的映射 $\varphi:\Sigma\to Y$。要区分「边界条件」与「边界理论（boundary theory）」：后者是一份数据。例如并不要求 $Z$ 是 $Y$ 的子流形；可以取一个映射 $Z\to Y$，并要求 $\partial\Sigma\to Y$ 提升为 $\partial\Sigma\to Z$。这是 sigma 模型（sigma model）中边界理论的典型特征。所谓 sigma 模型，就是某个场取值为到靶空间的映射的场论。

\srcpage{60}

把这一点用到有限群 Yang--Mills 的例子。此时靶是 $Y:=\mathrm{pt}/G$（看作群胚或叠）：一个映射 $\phi:\Sigma\to\mathrm{pt}/G$ 就是 $\Sigma$ 上的主 $G$-丛（principal $G$-bundle）。由于 $G$ 有限，它也就是 $\Sigma$ 上的 $G$-局部系统（local system）。一个映射 $Z\to Y$ 等价于一个带 $G$-作用的空间：总空间可以取主丛 $X\to Z$，也可以取 $Z\to\mathrm{pt}/G$ 与 $\mathrm{pt}\to\mathrm{pt}/G$ 的拉回。因此给定一个 $G$-空间 $X$，就得到边界数据。

$\Sigma$ 上的场空间，是 $X$ 上的 $G$-局部系统，连同 $\partial\Sigma\to\mathrm{pt}/G$ 沿 $X/G\to\mathrm{pt}/G$ 的提升。更几何地说，给定主 $G$-丛 $P\to\partial\Sigma$，可以造出配丛（associated bundle）
\[
P\times_G X:=\{(p,x)\mid (p\cdot g,x)\sim(p,g\cdot x)\}\longrightarrow\partial\Sigma .
\]
于是所谓一个场，就是 $\Sigma$ 上的主 $G$-丛 $P$，连同其配丛在 $\partial\Sigma$ 上的一个截面。等价地，这是从 $P|_{\partial\Sigma}$ 的总空间到 $X$ 的 $G$-等变映射。用物理的话说：我们在边界上把 sigma 模型与 $X$ 耦合起来，而体内是 $G$-规范理论（gauge theory）。

例如取 $X=G/H$。此时配丛的截面就是把结构群约化为主 $H$-丛。当 $H=\{1\}$ 时，这是主丛的一个平凡化；这就是 Dirichlet 边界条件，类似于要求一个函数在边界上调和。\footnote{还有 Neumann 边界条件，对应 $X=G/G$。}

在二维 TFT $Z_G$ 中，我们看到边界条件的范畴是 $\Rep G$。Dirichlet 边界条件对应正则表示，Neumann 边界条件对应另一个典范对象，即平凡表示。也可以考虑边界理论之间的界面（interface），它由表示之间的 $G$-等变映射给出。在 sigma 模型的例子里，边界数据 $X_1/G$ 与 $X_2/G$ 都带有到 $\mathrm{pt}/G$ 的映射。它们有典范的界面，即纤维积 $S:=X_1/G\times_{\mathrm{pt}/G}X_2/G$。在带边曲面上一部分边界提升到 $X_1/G$，另一部分提升到 $X_2/G$，在交叠处两者都要，于是映射到纤维积。

例如取 $X_1=G/H$、$X_2=G/K$，则 $X_1/G\cong\mathrm{pt}/H$、$X_2/G\cong\mathrm{pt}/K$。把 $\mathrm{pt}/H\to\mathrm{pt}/G$ 与 $\mathrm{pt}/K\to\mathrm{pt}/G$ 拉回，得到双陪集空间 $H\backslash G/K$。因此这两个边界条件之间的几何界面或 Lagrange 界面，由到双陪集空间的提升给出。$\mathrm{pt}/K$ 的自界面由到 $K\backslash G/K$ 的提升给出。于是我们的场是主 $G$-丛，边界的一部分约化为主 $H$-丛，另一部分约化为主 $K$-丛，而在交叠处有到双陪集空间的提升。也就是说，我们同时有一份 $H$-丛数据与一份 $K$-丛数据，以及它们诱导的 $G$-丛之间的同构。

Hecke 代数结构出现在复合界面的时候。设 $I_1$ 是 $X_1$ 与 $X_2$ 之间的界面，$I_2$ 是 $X_2$ 与 $X_3$ 之间的界面。物理上我们把用于 $X_2$ 的区域缩小；从足够远处看，它就像一个单一的缺陷（defect）。这种复合称为算子乘积展开（operator product expansion）。数学上，复合就是联系到 pair-of-chaps 协边的映射（TODO：图）。把它写在向量空间上，就得到映射
\[
\C[K\backslash G/K]\otimes\C[K\backslash G/K]\longrightarrow\C[K\backslash G/K],
\]
它正是 Hecke 代数的乘法。

例如当 $K=\{1\}$ 时 $K\backslash G/K=G$，Hecke 代数就是群代数。也就是说，Dirichlet 边界条件的自同态代数是 $\C[G]$。对有限群的表示论而言这是很强的结构，但在高维会非常有用。

\begin{remark}
存在一个版本，其中 $G$ 是任意紧 Lie 群，但一切都要导出化（derived）。我们的态空间不再是 $\mathrm{pt}/G$ 上的函数，而是
\[
H^*(\mathrm{pt}/G)=H^*_G(\mathrm{pt})=\C[\mathfrak h]^W,
\]
其中 $\mathfrak h$ 是 $G$ 的 Lie 代数的 Cartan 子代数，$W$ 是 Weyl 群。一个 Borel 子群会给出有趣的边界条件。
\end{remark}

然而，这些作为自同态代数出现的 Hecke 代数并不交换，例如 $\C[G/K]$ 这类对象作为表示的自同态代数。这对做谱理论是很大的障碍。但有时 Hecke 代数交换，例如 $H_{G\times G,G}=\End_{G\times G}\C[G]$。这是对任何群都可以同义反复地做的事，它总是交换的。

\srcpage{61}

把这个例子再挖深一点。主 $(G\times G)$-丛的数据就是两个主 $G$-丛，而约化到对角子群就是这两个丛之间的同构。所以我们要看的是：比方说在一个区间上的两个主丛，以及在边界点之外这两个丛之间的等同。把它想成「意大利饺」（ravioli）上的主 $G$-丛也许有帮助；这个非 Hausdorff 空间是 $I\amalg_{I\setminus\mathrm{pt}}I$。\yizhu{「意大利饺」指把两个区间沿各自内部挖去一点后的部分粘起来所得的图像；在二维场论中它对应圆上的 $G$-局部系统。}{B（nLab "ravioli"）}等价地，这是圆上的 $G$-局部系统：我们在 $S^1$ 的两个坐标卡上各粘一个 $G$-丛，于是需要在交叠处等同它们。也就是说，这个 Hecke 代数讲的是圆上的 $G$-理论，而不只是 $G\times G$ 的理论。

再换一种说法：我们在做 $G\times G$ 的 Yang--Mills 理论，看的是对角边界理论的自界面。$G\times G$ 的 Yang--Mills 的一个边界条件，等价于 $G$-理论的一个自界面。后一句更一般：$G\times H$ 理论的边界条件，等价于从 $G$-理论到 $H$-理论的界面。一个 $\C[G\times G]$-模等价于一个 $\C[G]\otimes\C[G]$-模，等价于一个 $\C[G]^{\mathrm{op}}\otimes\C[G]$-模，也等价于一个 $\C[G]$-双模。

这是场论的一般特点：一旦理解了边界条件，就理解了界面，反之亦然。

于是回到一个重要问题：为什么 Hecke 代数 $H_{G\times G,G}=\C[G/G]$ 交换？从 Hecke 的角度看，它没有理由交换。但从场论的角度看有理由：$H_{G\times G,G}$ 是 $Z_G$ 的平凡自界面代数。数学上，它是 $\Rep G$ 的恒等函子的自同态代数。对任何范畴 $\cC$，$\End(\id_\cC)$ 称为 $\cC$ 的 Bernstein 中心，记作 $Z(\cC)$。

交换性正来自这里。我们有两个不同的复合方向，可以用它把 $A\circ B$ 旋转成 $B\circ A$；其他 Hecke 代数没有这个性质。这与「2 阶及以上的同伦群交换」的证明密切相关。这类论证有时称为 Eckmann--Hilton 论证：Bernstein 中心有两个相容的幺半结构，即复合与张量积，它们彼此交换。\yizhu{Eckmann--Hilton 论证是代数拓扑中的标准技巧：两个相互分配的幺半结构必重合且交换，由此可证 $\pi_n$（$n\ge2$）交换。}{B（nLab "Eckmann-Hilton argument"）}一个标准而不难的练习是：带有两个相容幺半结构的集合是交换幺半群，两个乘法重合且交换。另一个相关论证说明：若 $G$ 是拓扑群，则 $\pi_1(G)$ 交换——我们有两个乘法，一个来自 $G$，一个来自环路复合，它们交换。

设 $A$ 是结合代数；则 Bernstein 中心 $Z(\mathrm{Mod}_A)=\End(A)$，这里把 $A$ 看作 $(A,A)$-双模。这是因为 $\mathrm{Mod}_A$ 上的恒等函子等价于与 $(A,A)$-双模 $A$ 作张量积。作为 $(A,A)$-双模，$A$ 的自同态是与左右乘法都交换的东西，这正是 $A$ 的中心。这就是 Bernstein 中心被称为「中心」的原因：它是代数中心的推广。

于是对任何范畴 $\cC$ 我们都有交换代数 $Z(\cC)$，所以应当对它作谱分解。令 $\cM_\cC:=\Spec(Z\cC)$；从这个空间能学到什么？例如对 $\cC=\mathrm{Mod}_A$，它是 $\Spec(Z(A))$。这并不总是有用：有时 $Z(A)=0$，例如微分算子的 Weyl 代数。

一般来说，高维 QFT 会以类似方式带有内蕴交换的对象，它们来自局部算子与缺陷。在 QFT 中，局部算子（local operator）是某种测量或可观测量，定域在时空 $M$ 中一个点的附近。可以取该点周围半径为 $\varepsilon$ 的球面，并把球从时空中挖掉，于是得到一个从 $S^{n-1}_\varepsilon$ 到空集的协边。这样我们就得到 $S^{n-1}_\varepsilon$ 的态空间上的泛函。然后令 $\varepsilon\to0$；这允许场在我们测量的那个点上奇异，因为 $\varepsilon$ 永远不等于 $0$。这个算子 $\cO$ 的期望值是对场空间的路径积分：
\[
\langle \cO(x)\rangle=\int \cO_x(\varphi)e^{-S(\varphi)}\,D\varphi .
\]
（启发式地）我们在 $M\setminus x$ 上的场空间上积分。所以在物理中，局部算子就是给出期望值以及关联函数的东西。后者指在不同点放多个局部算子，并把它们都插入路径积分。

在拓扑场论（更一般地，共形场论）中，不必担心球面的大小 $\varepsilon$，也就不必担心极限。球面只对应一个向量空间，它就是局部算子空间。这称为态-算子对应（state-operator correspondence）。算子的期望值由挖去一个小球面决定：剩下的流形是从 $S^{n-1}$ 到 $\varnothing$ 的协边，于是定义了从局部算子到 $\C$ 的函数。这就是期望值映射。

\srcpage{62}

具体地，对二维拓扑 Yang--Mills 理论（$G$ 有限）做这件事，局部算子是 $Z_G(S^1)=\C[G/G]$。于是 $G/G$ 上的 $\delta$-函数构成一组基；每个 $\delta$-函数指定主 $G$-丛绕一个环的可能单值化（monodromy）。于是局部算子的期望值可以解释为：数 $M\setminus x$ 上、绕 $x$ 有指定单值化的 $G$-局部系统的个数。

关键是：$H_{G\times G,G}$ 是 $Z_G$ 中局部算子的代数，而算子乘积展开（两个点碰撞）就是裤子协边或环协边。这些说法表达同一件事：局部算子空间上有代数结构，而由于拓扑原因这个代数交换。交换就意味着可以作谱分解。数论中 Hecke 代数的交换性类似：从 Hecke 代数的角度看，它们没有理由交换；但我们会把它们解释为局部算子代数，从而说明其交换性。
